%% kapa: exam=1 round=32 year=2021 period=1 form=A subject=economics subject_ko=경제학원론 seq=29 number=69
%% source: Q-Net 감정평가사 기출문제 (cst00309) artlSeq=5208397 file=extracted/1008/5208397/A형/2021년도 제32회 감정평가사 1차 1교시 A형.pdf
%% conversion: kordoc 4.24.51 -> tools/gen1.py (auto); review=proofread-sonnet; flags=none
\begin{question}{69}
표의 기존 가정에 따라 독립투자승수를 계산했다. 계산된 승수를 하락시키는 가정의 변화를 모두 고른 것은?
\par{\small\centering
\begin{tabular}{|p{0.36\linewidth}|c|p{0.36\linewidth}|}\hline
\multicolumn{2}{|c|}{기존 가정} & \multicolumn{1}{c|}{가정의 변화} \\ \hline
ㄱ. 생산자들은 고정된 가격에 추가적인 생산물을 공급한다. & $\to$ & 총공급 곡선이 수직이다. \\ \hline
ㄴ. 이자율은 고정이다. & $\to$ & 이자율 상승에 따라 투자가 감소한다. \\ \hline
ㄷ. 정부지출과 세금은 없다. & $\to$ & 정부지출과 세금이 모두 외생적으로 증가한다. \\ \hline
ㄹ. 수출과 수입은 모두 영(0)이다. & $\to$ & 수출과 수입이 모두 외생적으로 증가한다. \\ \hline
\end{tabular}\par}
\begin{choices}
  \item ㄱ, ㄴ
  \item ㄱ, ㄷ
  \item ㄴ, ㄷ
  \item ㄴ, ㄹ
  \item ㄷ, ㄹ
\end{choices}
\end{question}
\answer{1}
